%% Invariant Normalization Formalization for ICML 2026
%% This file contains LaTeX content for the methods section
%% describing the semantic-preserving invariant normalization procedure.

%% Required packages (add to your main document preamble):
%% \usepackage{algorithm}
%% \usepackage{algorithmic}
%% \usepackage{amsmath}
%% \usepackage{amssymb}
%% \usepackage{booktabs}
%% \usepackage{tikz}
%% \usetikzlibrary{shapes.geometric, arrows.meta, positioning}
%% \usepackage{enumitem}  % For nosep option in itemize
%% \usepackage{listings}  % For code listings
%% \usepackage{xcolor}    % For colorbox highlighting

%% ============================================================================
%% SECTION: FORMAL DEFINITIONS
%% ============================================================================

%% Use this in your Methods/Data Pipeline section:

\subsection{Invariant Normalization}
\label{sec:invariant-normalization}

Automated verification tools (e.g., UAutomizer~\cite{heizmann2013ultimate}, CPAchecker, SeaHorn) produce invariants containing syntactic artifacts: type casts, tautological clauses (e.g., \texttt{x <= x}), and excessive parentheses. We apply semantic-preserving normalization to obtain canonical representations. While we evaluate on UAutomizer invariants, the normalization rules are general and apply to any verifier producing C-like boolean expressions.

%% ============================================================================
%% RUNNING EXAMPLE (demonstrates all normalization issues)
%% ============================================================================

\begin{figure}[t]
\centering
\small
\begin{tabular}{@{}p{0.95\linewidth}@{}}
\toprule
\textbf{Program:} Computing sum and sum of squares \\
\midrule
\begin{verbatim}
long long sum = 0, sqsum = 0;
int c = 0;
while (c < n) {
  INVARIANT();
  c++; sum += c; sqsum += c * c;
}
assert(2*sum == c*(c+1) && 6*sqsum == c*(c+1)*(2*c+1));
\end{verbatim} \\
\midrule
\textbf{Raw invariant from UAutomizer} (174 characters): \\
\midrule
\texttt{(((\colorbox{pink}{(\_\_int128)} sqsum * 6) == (((\colorbox{pink}{(\_\_int128)} (\colorbox{pink}{(long long)} c * c) * 3) + c)} \\
\texttt{+ (2 * (\colorbox{pink}{(\_\_int128)} (\colorbox{pink}{(long long)} c * c) * c)))) \&\& ((\colorbox{pink}{(\_\_int128)} 2 * sum)} \\
\texttt{== ((\colorbox{pink}{(long long)} c * c) + c)))} \\
\midrule
\textbf{After normalization} (66 characters, \textbf{62\% reduction}): \\
\midrule
\texttt{sqsum * 6 == c * c * 3 + c + 2 * c * c * c \&\& 2 * sum == c * c + c} \\
\bottomrule
\end{tabular}
\caption{\textbf{Running example} of invariant normalization. The raw invariant contains \colorbox{pink}{7 type casts} (highlighted) and 21 pairs of parentheses. After normalization: (1) integral casts are removed, (2) parentheses are minimized using operator precedence. The result is 62\% shorter and semantically equivalent.}
\label{fig:running-example}
\end{figure}

%% Example with tautology (brief inline mention)
% Another common pattern involves tautological clauses. For instance, the invariant:
% \texttt{((2*sum == i*i - i) \&\& (i <= n)) \&\& \colorbox{yellow}{(n <= n)}}
% normalizes to simply \texttt{2*sum == i*i - i \&\& i <= n}, removing the
% trivially true clause \texttt{n <= n} (65\% reduction).

\begin{definition}[Invariant Expression]
\label{def:invariant-expr}
An \emph{invariant expression} $\varphi$ over program variables $V$ is defined by the grammar:
\begin{align*}
\varphi &::= \mathit{true} \mid \mathit{false} \mid \varphi_1 \land \varphi_2 \mid \varphi_1 \lor \varphi_2 \mid \neg \varphi \mid e_1 \bowtie e_2 \\
e &::= x \mid c \mid e_1 \oplus e_2 \mid {-}e \mid (T)\, e
\end{align*}
where $x \in V$ is a program variable, $c \in \mathbb{Z}$ is an integer constant, $\bowtie \in \{<, \leq, =, \neq, \geq, >\}$ is a comparison operator, $\oplus \in \{+, -, \times, /, \%\}$ is an arithmetic operator, and $(T)\, e$ denotes a type cast to integral type $T$.
\end{definition}

\begin{definition}[Semantic Equivalence]
\label{def:semantic-equiv}
Two invariant expressions $\varphi_1$ and $\varphi_2$ are \emph{semantically equivalent}, written $\varphi_1 \equiv \varphi_2$, if for all valuations $\sigma: V \to \mathbb{Z}$, we have $\llbracket \varphi_1 \rrbracket_\sigma = \llbracket \varphi_2 \rrbracket_\sigma$, where $\llbracket \cdot \rrbracket_\sigma$ denotes the standard interpretation of expressions under $\sigma$.
\end{definition}

%% ============================================================================
%% SECTION: REWRITE RULES (as a figure)
%% ============================================================================

\begin{figure}[t]
\centering
\small
\begin{tabular}{@{}ll@{}}
\toprule
\textbf{Rule} & \textbf{Transformation} \\
\midrule
\multicolumn{2}{@{}l}{\textit{Cast Elimination} (integral types only)} \\
\textsc{Cast} & $(T)\, e \;\longrightarrow\; e$ \quad where $T \in \{\texttt{int}, \texttt{long}, \texttt{long long}, \ldots\}$ \\
\midrule
\multicolumn{2}{@{}l}{\textit{Tautology Elimination}} \\
\textsc{TautConj-L} & $\mathit{true} \land \varphi \;\longrightarrow\; \varphi$ \\
\textsc{TautConj-R} & $\varphi \land \mathit{true} \;\longrightarrow\; \varphi$ \\
\textsc{TautRefl-$\leq$} & $\varphi \land (e \leq e) \;\longrightarrow\; \varphi$ \\
\textsc{TautRefl-$\geq$} & $\varphi \land (e \geq e) \;\longrightarrow\; \varphi$ \\
\textsc{TautRefl-$=$} & $\varphi \land (e = e) \;\longrightarrow\; \varphi$ \\
\textsc{TautConst} & $\varphi \land (c_1 \bowtie c_2) \;\longrightarrow\; \varphi$ \quad\text{if } $c_1 \bowtie c_2$ is true \\
\midrule
\multicolumn{2}{@{}l}{\textit{Contradiction Propagation}} \\
\textsc{ContraConj-L} & $\mathit{false} \land \varphi \;\longrightarrow\; \mathit{false}$ \\
\textsc{ContraConj-R} & $\varphi \land \mathit{false} \;\longrightarrow\; \mathit{false}$ \\
\textsc{ContraRefl-$<$} & $\varphi \land (e < e) \;\longrightarrow\; \mathit{false}$ \\
\textsc{ContraRefl-$>$} & $\varphi \land (e > e) \;\longrightarrow\; \mathit{false}$ \\
\textsc{ContraRefl-$\neq$} & $\varphi \land (e \neq e) \;\longrightarrow\; \mathit{false}$ \\
\textsc{ContraDisj-L} & $\mathit{false} \lor \varphi \;\longrightarrow\; \varphi$ \\
\textsc{ContraDisj-R} & $\varphi \lor \mathit{false} \;\longrightarrow\; \varphi$ \\
\textsc{TautDisj} & $\mathit{true} \lor \varphi \;\longrightarrow\; \mathit{true}$ \\
\bottomrule
\end{tabular}
\caption{Semantic-preserving rewrite rules for invariant normalization. Tautology rules (69.8\% of invariants affected) and contradiction rules (included for completeness; rare in practice) together handle \emph{trivially decidable clauses}. Rules are applied exhaustively in a single bottom-up pass over the AST.}
\label{fig:rewrite-rules}
\end{figure}

%% ============================================================================
%% SECTION: ALGORITHM
%% ============================================================================

\begin{algorithm}[t]
\caption{Invariant Normalization}
\label{alg:normalize}
\begin{algorithmic}[1]
\REQUIRE Raw invariant expression $\varphi$ (as AST)
\ENSURE Normalized expression $\varphi'$ such that $\varphi' \equiv \varphi$
\STATE $\varphi' \gets \textsc{Transform}(\varphi)$
\STATE \textbf{return} \textsc{PrettyPrint}($\varphi'$) \COMMENT{Minimal parentheses}
\end{algorithmic}

\vspace{0.5em}
\begin{algorithmic}[1]
\FUNCTION{\textsc{Transform}}{$\varphi$}
    \IF{$\varphi = (T)\, e$ \AND $T$ is integral}
        \STATE \textbf{return} \textsc{Transform}($e$)
    \ELSIF{$\varphi = \varphi_1 \land \varphi_2$}
        \STATE $\varphi_1' \gets \textsc{Transform}(\varphi_1)$
        \STATE $\varphi_2' \gets \textsc{Transform}(\varphi_2)$
        \IF{\textsc{IsTautology}($\varphi_1'$)}
            \STATE \textbf{return} $\varphi_2'$
        \ELSIF{\textsc{IsTautology}($\varphi_2'$)}
            \STATE \textbf{return} $\varphi_1'$
        \ELSIF{\textsc{IsContradiction}($\varphi_1'$) \OR \textsc{IsContradiction}($\varphi_2'$)}
            \STATE \textbf{return} $\mathit{false}$
        \ENDIF
        \STATE \textbf{return} $\varphi_1' \land \varphi_2'$
    \ELSIF{$\varphi = \varphi_1 \lor \varphi_2$}
        \STATE $\varphi_1' \gets \textsc{Transform}(\varphi_1)$
        \STATE $\varphi_2' \gets \textsc{Transform}(\varphi_2)$
        \IF{\textsc{IsTautology}($\varphi_1'$) \OR \textsc{IsTautology}($\varphi_2'$)}
            \STATE \textbf{return} $\mathit{true}$
        \ELSIF{\textsc{IsContradiction}($\varphi_1'$)}
            \STATE \textbf{return} $\varphi_2'$
        \ELSIF{\textsc{IsContradiction}($\varphi_2'$)}
            \STATE \textbf{return} $\varphi_1'$
        \ENDIF
        \STATE \textbf{return} $\varphi_1' \lor \varphi_2'$
    \ELSE
        \STATE \textbf{return} $\varphi$ \COMMENT{Atoms unchanged}
    \ENDIF
\ENDFUNCTION
\end{algorithmic}

\vspace{0.5em}
\begin{algorithmic}[1]
\FUNCTION{\textsc{IsTautology}}{$\varphi$}
    \IF{$\varphi = \mathit{true}$ \OR $\varphi = 1$}
        \STATE \textbf{return} \TRUE
    \ELSIF{$\varphi = (e_1 \bowtie e_2)$ \AND $e_1 \equiv_{\text{syn}} e_2$ \AND $\bowtie \in \{\leq, \geq, =\}$}
        \STATE \textbf{return} \TRUE
    \ELSIF{$\varphi = (c_1 \bowtie c_2)$ \AND $c_1 \bowtie c_2$ evaluates to true}
        \STATE \textbf{return} \TRUE
    \ENDIF
    \STATE \textbf{return} \FALSE
\ENDFUNCTION
\end{algorithmic}

\vspace{0.5em}
\begin{algorithmic}[1]
\FUNCTION{\textsc{IsContradiction}}{$\varphi$}
    \IF{$\varphi = \mathit{false}$ \OR $\varphi = 0$}
        \STATE \textbf{return} \TRUE
    \ELSIF{$\varphi = (e_1 \bowtie e_2)$ \AND $e_1 \equiv_{\text{syn}} e_2$ \AND $\bowtie \in \{<, >, \neq\}$}
        \STATE \textbf{return} \TRUE
    \ELSIF{$\varphi = (c_1 \bowtie c_2)$ \AND $c_1 \bowtie c_2$ evaluates to false}
        \STATE \textbf{return} \TRUE
    \ENDIF
    \STATE \textbf{return} \FALSE
\ENDFUNCTION
\end{algorithmic}
\end{algorithm}

%% ============================================================================
%% SECTION: SOUNDNESS
%% ============================================================================

\begin{proposition}[Soundness]
\label{prop:soundness}
For any invariant expression $\varphi$, the normalized expression $\varphi' = \textsc{Normalize}(\varphi)$ satisfies $\varphi' \equiv \varphi$ under mathematical integer semantics.
\end{proposition}

\begin{proof}[Proof sketch]
Tautology and contradiction elimination follow directly from Boolean algebra. Cast elimination is sound under SV-COMP's standard \emph{mathematical integer semantics} (infinite domain $\mathbb{Z}$, no overflow), where integral type casts serve only for type-checking, not overflow protection. (Pointer casts are preserved but absent in our ReachSafety dataset.)

Empirical validation: 0\% of 13{,}776 normalized invariants failed re-verification.
\end{proof}

%% ============================================================================
%% SECTION: EXAMPLE (as a figure)
%% ============================================================================

\begin{figure}[t]
\centering
\small
\begin{tabular}{@{}p{0.95\linewidth}@{}}
\toprule
\textbf{Raw invariant from UAutomizer:} \\
\midrule
\texttt{((((prime\_count <= x) \&\& ((((long long) prime\_count + 1) \% 4294967296) <= x))} \\
\texttt{\&\& ((((long long) 2 + prime\_count) \% 4294967296) <= ((long long) n + 1)))} \\
\texttt{|| ((prime\_count == 0) \&\& (n == 0)))} \\
\midrule
\textbf{After normalization:} \\
\midrule
\texttt{prime\_count <= x \&\& (prime\_count + 1) \% 4294967296 <= x} \\
\texttt{\&\& (2 + prime\_count) \% 4294967296 <= n + 1} \\
\texttt{|| prime\_count == 0 \&\& n == 0} \\
\bottomrule
\end{tabular}
\caption{Example of invariant normalization. The raw invariant contains 7 type casts (\texttt{(long long)}) and excessive parentheses. After normalization, casts are removed and parentheses are minimized using operator precedence, reducing the expression length by 33\% while preserving semantic equivalence. In our training set of 4{,}000 invariants, 44.8\% contain type casts and 69.8\% contain tautological clauses, with an average length reduction of 32.8\% after normalization.}
\label{fig:normalization-example}
\end{figure}

%% ============================================================================
%% SECTION: TAUTOLOGY REMOVAL EXAMPLE
%% ============================================================================

\begin{figure}[t]
\centering
\small
\begin{tabular}{@{}p{0.95\linewidth}@{}}
\toprule
\textbf{Invariant with tautological clauses:} \\
\midrule
\texttt{(x <= x) \&\& (y > 0) \&\& (0 < 1) \&\& (z >= z)} \\
\midrule
\textbf{After tautology elimination:} \\
\midrule
\texttt{y > 0} \\
\bottomrule
\end{tabular}
\caption{Example of tautology elimination. The clauses \texttt{x <= x}, \texttt{0 < 1}, and \texttt{z >= z} are trivially true and removed from the conjunction, leaving only the semantically meaningful constraint.}
\label{fig:tautology-example}
\end{figure}

%% ============================================================================
%% SECTION: CONTRADICTION ELIMINATION EXAMPLE (real example from 9097_1.c)
%% ============================================================================

\begin{figure}[t]
\centering
\small
\begin{tabular}{@{}p{0.95\linewidth}@{}}
\toprule
\textbf{Program (simplified):} \\
\midrule
\begin{verbatim}
unsigned int i = 0, sum = 0;
while (i < 10) {
  INVARIANT_MARKER_1();
  if (i % 2 == 0) sum += i; else sum -= i;
  i++;
}
assert(sum == -5);
\end{verbatim} \\
\midrule
\textbf{Raw invariant from UAutomizer (454 chars):} \\
\midrule
\texttt{((((((((((((i == 4) \&\& (((long long) 2 + sum) == 0)) ||} \\
\texttt{~~((((0 <= ((long long) 5 + sum)) \&\& ...} \\
\texttt{~~~~\&\& (i == 10)) \&\& \colorbox{yellow}{(1 <= 0)})) || ...} \\
\midrule
\textbf{After normalization (232 chars, 49\% reduction):} \\
\midrule
\texttt{i == 4 \&\& 2 + sum == 0 || 3 == i \&\& sum == 1 ||} \\
\texttt{sum == 0 \&\& i == 1 || 3 + sum == 0 \&\& 6 == i || ...} \\
\bottomrule
\end{tabular}
\caption{Real example showing contradiction elimination. The raw invariant contains the clause \texttt{(1 <= 0)}, which is always false. By the rule $\varphi \land \mathit{false} \to \mathit{false}$, the entire disjunct containing it is eliminated. Combined with cast removal and parenthesis minimization, the invariant length is reduced by 49\%.}
\label{fig:contradiction-example}
\end{figure}

%% ============================================================================
%% SECTION: LLM-BASED SIMPLIFICATION AND WEAKENING
%% ============================================================================

\subsection{LLM-Based Invariant Simplification}
\label{sec:llm-simplification}

While AST-based normalization removes syntactic noise, many invariants remain verbose due to \emph{semantic} complexity: enumerated cases, redundant bounds implied by program context, or overly strong constraints that are not necessary to prove the target property. These patterns require reasoning beyond local syntactic rewrites. We employ an LLM to perform \emph{invariant weakening}---producing simpler invariants that are logically weaker but still sufficient for verification.

\begin{definition}[Invariant Weakening]
\label{def:weakening}
An invariant $\varphi'$ is a \emph{valid weakening} of $\varphi$ at loop location $l$ with respect to program $P$ and target property $\psi$ if:
\begin{enumerate}[nosep,leftmargin=*]
    \item \textbf{(Weaker or equivalent)} $\varphi \models \varphi'$, i.e., $\varphi'$ is logically implied by $\varphi$
    \item \textbf{(Inductive)} $\varphi'$ is a valid loop invariant at $l$ in $P$
    \item \textbf{(Sufficient)} $\varphi'$ is strong enough to prove the target property $\psi$
\end{enumerate}
\end{definition}

The key insight is that automated verifiers often produce invariants stronger than necessary---they capture precise execution states rather than the minimal constraints needed for correctness. An LLM can recognize patterns and generalize, but may produce incorrect candidates. We thus employ a \emph{generate-and-verify} paradigm with formal verification in the loop.

\paragraph{Simplification Patterns.}
The LLM is prompted to apply semantic simplifications that complement AST-based normalization:
\begin{itemize}[nosep,leftmargin=*]
    \item \textbf{Range generalization}: $\bigvee_{i \in \{1,2,3,4,5\}} (x = i) \;\rightsquigarrow\; 1 \leq x \land x \leq 5$
    \item \textbf{Bound extraction}: Extract common constraints from disjuncts
    \item \textbf{Constant elimination}: Remove constraints on variables that are constant throughout loop execution
    \item \textbf{Redundancy removal}: Eliminate constraints implied by program semantics (e.g., $x \geq 0$ when $x$ is initialized to $0$ and only incremented)
\end{itemize}

\paragraph{Best-of-N Sampling with Verification.}
For each verbose invariant $\varphi$ (defined as $|\varphi| \geq \tau$ characters after normalization, where $\tau = 20$ in our experiments), we generate $N$ candidate simplifications $\{\varphi'_1, \ldots, \varphi'_N\}$ using temperature sampling ($T=1.0$) from a reasoning-capable LLM (Kimi-K2~\cite{team2025kimik2}). Each candidate is then verified in parallel using the formal verifier (UAutomizer).

\begin{definition}[Quality Grade]
\label{def:quality-grade}
Given a candidate $\varphi'$, program $P$, target property $\psi$, and baseline verification time $t_{\text{base}}$, we define the quality grade $Q(\varphi') \in \{0, 1, 2, 3\}$:
\begin{align*}
Q(\varphi') = \begin{cases}
0 & \text{if } \varphi' \text{ is not a valid invariant (verification fails)} \\
1 & \text{if } \varphi' \text{ is valid but } \psi \text{ cannot be proved using } \varphi' \\
2 & \text{if } \varphi' \text{ proves } \psi \text{ but } t(\varphi') \geq t_{\text{base}} \\
3 & \text{if } \varphi' \text{ proves } \psi \text{ and } t(\varphi') < t_{\text{base}} \quad \text{(speedup)}
\end{cases}
\end{align*}
where $t(\varphi')$ is the verification time using $\varphi'$ as an assumption.
\end{definition}

\begin{algorithm}[t]
\caption{LLM-Based Invariant Simplification}
\label{alg:llm-simplify}
\begin{algorithmic}[1]
\REQUIRE Normalized invariant $\varphi$, program $P$, target $\psi$, baseline time $t_{\text{base}}$, LLM $\mathcal{M}$, $N$ candidates
\ENSURE Simplified invariant $\varphi^*$ with $Q(\varphi^*) > 0$, or $\bot$ if no valid simplification
\IF{$|\varphi| < \tau$}
    \STATE \textbf{return} $(\varphi, \textsc{Verify}(\varphi, P, \psi, t_{\text{base}}))$ \COMMENT{Non-verbose: verify as-is}
\ENDIF
\STATE $\mathcal{C} \gets \textsc{LLM-Generate}(\mathcal{M}, P, \varphi, N)$ \COMMENT{Sample $N$ candidates}
\STATE $\mathcal{C} \gets \textsc{Deduplicate}(\mathcal{C})$
\STATE $\mathcal{R} \gets \varnothing$
\FORALL{$\varphi' \in \mathcal{C}$ \textbf{in parallel}}
    \STATE $q \gets \textsc{Verify}(\varphi', P, \psi, t_{\text{base}})$
    \IF{$q > 0$}
        \STATE $\mathcal{R} \gets \mathcal{R} \cup \{(\varphi', q)\}$
    \ENDIF
\ENDFOR
\IF{$\mathcal{R} \neq \varnothing$}
    \STATE \textbf{return} $\arg\max_{(\varphi', q) \in \mathcal{R}} q$ \COMMENT{Best quality candidate}
\ELSE
    \STATE $q_{\text{fallback}} \gets \textsc{Verify}(\varphi, P, \psi, t_{\text{base}})$ \COMMENT{Fallback to normalized}
    \STATE \textbf{return} $(\varphi, q_{\text{fallback}})$ if $q_{\text{fallback}} > 0$ else $\bot$
\ENDIF
\end{algorithmic}
\end{algorithm}

\begin{figure}[t]
\centering
\small
\begin{tabular}{@{}p{0.95\linewidth}@{}}
\toprule
\textbf{Program:} Binary search tracking \\
\midrule
\begin{lstlisting}[language=C, basicstyle=\ttfamily\footnotesize, xleftmargin=1em]
int target = 10, size = 5, res = 0, i = 0;
while (i < size) {
  INVARIANT_MARKER_1();
  if (arr[i] == target) res = i + 1;
  i++;
}
assert(res >= 0);
\end{lstlisting} \\
\midrule
\textbf{Normalized invariant} (137 characters): \\
\midrule
\texttt{target == 10 \&\& 5 == size \&\& res == 0 \&\& i == 0 ||} \\
\texttt{target == 10 \&\& 5 == size \&\& i == 1 \&\& res <= 1 ||} \\
\texttt{target == 10 \&\& 5 == size \&\& 2 == i \&\& res <= 3 || 5 == size \&\& 3 < i} \\
\midrule
\textbf{After LLM simplification} (18 characters, \textbf{87\% reduction}): \\
\midrule
\texttt{i >= 0 \&\& res >= 0} \\
\midrule
\textbf{Rationale}: Simplified to basic bounds. \texttt{size == 5} and \texttt{target == 10} are constants initialized before the loop and never modified. The enumerated cases for \texttt{i} generalize to \texttt{i >= 0}. \\
\bottomrule
\end{tabular}
\caption{\textbf{LLM-based simplification example}. The normalized invariant enumerates specific states; the LLM recognizes that (1)~\texttt{target} and \texttt{size} are loop-invariant constants, (2)~the enumerated \texttt{i} values generalize to a lower bound, and (3)~the constraint \texttt{res >= 0} suffices to prove the target. The simplified invariant is verified to be inductive and sufficient.}
\label{fig:llm-simplification-example}
\end{figure}

\begin{figure}[t]
\centering
\small
\begin{tabular}{@{}p{0.95\linewidth}@{}}
\toprule
\textbf{Normalized invariant} (disjunctive form, 156 characters): \\
\midrule
\texttt{b == 1 \&\& 5 == i || b == 1 \&\& 2 == i || b == 1 \&\& i == 4 ||} \\
\texttt{b == 1 \&\& 3 == i || b == 1 \&\& i == 1 || b == 1 \&\& 6 == i} \\
\midrule
\textbf{After LLM simplification} (14 characters): \\
\midrule
\texttt{1 <= i \&\& i <= 6} \\
\midrule
\textbf{Rationale}: The enumerated values $i \in \{1,2,3,4,5,6\}$ form a contiguous range. The constraint \texttt{b == 1} is redundant (always true in this context). \\
\bottomrule
\end{tabular}
\caption{\textbf{Range generalization example}. The LLM recognizes that six disjuncts enumerate consecutive values and generalizes to an interval constraint. This pattern is common in verifier output for bounded loops.}
\label{fig:range-generalization-example}
\end{figure}

\paragraph{Soundness via Verification.}
Unlike AST-based normalization, LLM-generated simplifications are not guaranteed to be correct. Soundness is ensured by formal verification: each candidate must pass both (1)~\emph{correctness checking}---verifying that $\varphi'$ holds at the loop location, and (2)~\emph{usefulness checking}---verifying that assuming $\varphi'$ is sufficient to prove the target property $\psi$. Only candidates with $Q(\varphi') > 0$ are retained for training.

\begin{proposition}[Pipeline Soundness]
\label{prop:pipeline-soundness}
Let $\varphi^*$ be the output of Algorithm~\ref{alg:llm-simplify}. If $\varphi^* \neq \bot$, then $\varphi^*$ is a valid loop invariant for program $P$ that is sufficient to prove target property $\psi$.
\end{proposition}

\begin{proof}
By construction, $\varphi^*$ is returned only if $\textsc{Verify}(\varphi^*, P, \psi, t_{\text{base}}) > 0$, which requires passing the formal verifier's correctness and usefulness checks.
\end{proof}

%% ============================================================================
%% COMPACT VERSION FOR SPACE-CONSTRAINED PAPERS
%% ============================================================================

%% If you need a more compact version for the main paper (with details in appendix),
%% use this paragraph instead:

% \paragraph{Invariant Normalization.}
% Raw invariants from verification tools contain syntactic noise: type casts 
% (e.g., \texttt{(long long) x}), tautological clauses (e.g., $x \leq x$), and 
% excessive parentheses. We apply semantic-preserving rewrite rules to obtain 
% canonical forms: (1)~integral type casts are removed, (2)~tautologies are 
% eliminated from conjunctions (and contradictions from disjunctions), and 
% (3)~parentheses are minimized using C operator precedence. See 
% Appendix~\ref{app:normalization} for formal definitions and the complete 
% algorithm. This normalization reduces average invariant length by 32.8\% in 
% our dataset while preserving semantic equivalence.
%
% \paragraph{LLM-Based Simplification.}
% For verbose invariants (>20 characters), we employ an LLM to perform semantic
% simplification: generalizing enumerated cases to ranges, removing redundant 
% constraints implied by program context, and producing weaker-but-sufficient 
% invariants. We use Best-of-N sampling with formal verification in the loop 
% to ensure soundness. See Appendix~\ref{app:llm-simplification} for details.

%% ============================================================================
%% COMBINED PIPELINE FIGURE
%% ============================================================================

\begin{figure}[t]
\centering
\small
\begin{tikzpicture}[
    node distance=0.8cm and 1.2cm,
    box/.style={rectangle, draw, rounded corners, minimum width=2.2cm, minimum height=0.7cm, align=center, font=\scriptsize},
    stage/.style={box, fill=blue!10},
    decision/.style={diamond, draw, aspect=2, inner sep=1pt, font=\scriptsize},
    arrow/.style={->, >=stealth, thick}
]
% Stage 1: AST Normalization
\node[stage] (raw) {Raw Invariant\\$\varphi_{\text{raw}}$};
\node[stage, right=of raw] (parse) {Parse to AST};
\node[stage, right=of parse] (rewrite) {Apply Rewrite\\Rules};
\node[stage, right=of rewrite] (pretty) {Pretty Print\\$\varphi_{\text{norm}}$};

% Stage 2: LLM Simplification
\node[decision, below=1.2cm of pretty] (verbose) {Verbose?\\$|\varphi| \geq \tau$};
\node[stage, below left=0.8cm and 0.3cm of verbose] (llm) {LLM Generate\\$N$ Candidates};
\node[stage, below=0.8cm of llm] (verify) {Parallel\\Verification};
\node[decision, below=0.8cm of verify] (pass) {$Q > 0$?};
\node[stage, below right=0.8cm and 0.3cm of verbose] (direct) {Verify\\As-Is};
\node[stage, below=2.5cm of verbose] (output) {$\varphi^*$};

% Arrows - Stage 1
\draw[arrow] (raw) -- (parse);
\draw[arrow] (parse) -- (rewrite);
\draw[arrow] (rewrite) -- (pretty);

% Arrows - Stage 2
\draw[arrow] (pretty) -- (verbose);
\draw[arrow] (verbose) -- node[left, font=\tiny] {Yes} (llm);
\draw[arrow] (verbose) -- node[right, font=\tiny] {No} (direct);
\draw[arrow] (llm) -- (verify);
\draw[arrow] (verify) -- (pass);
\draw[arrow] (pass) -- node[left, font=\tiny] {Yes} (output);
\draw[arrow] (pass.east) -- ++(0.5,0) |- node[right, font=\tiny, pos=0.25] {No} (llm.east);
\draw[arrow] (direct) |- (output);

% Stage labels
\node[above=0.1cm of parse, font=\scriptsize\bfseries] {Stage 1: AST Normalization};
\node[below=0.1cm of verbose, xshift=-1.5cm, font=\scriptsize\bfseries] {Stage 2: LLM Simplification};
\end{tikzpicture}
\caption{\textbf{Two-stage invariant curation pipeline}. Stage~1 applies semantic-preserving AST rewrites (cast elimination, tautology removal). Stage~2 uses an LLM with verification-in-the-loop to further simplify verbose invariants via semantic weakening.}
\label{fig:pipeline}
\end{figure}

%% ============================================================================
%% COMBINED EXAMPLE SHOWING BOTH STAGES
%% ============================================================================

\begin{figure*}[t]
\centering
\small
\begin{tabular}{@{}p{0.95\linewidth}@{}}
\toprule
\textbf{Stage 0: Raw invariant from UAutomizer} (289 characters): \\
\midrule
\texttt{(((\colorbox{pink!50}{(\_\_int128)} b == 1) \&\& (\colorbox{pink!50}{(\_\_int128)} 5 == i)) || ((\colorbox{pink!50}{(\_\_int128)} b == 1) \&\& (\colorbox{pink!50}{(\_\_int128)} 2 == i))} \\
\texttt{|| ((\colorbox{pink!50}{(\_\_int128)} b == 1) \&\& (\colorbox{pink!50}{(\_\_int128)} i == 4)) || ((\colorbox{pink!50}{(\_\_int128)} b == 1) \&\& (\colorbox{pink!50}{(\_\_int128)} 3 == i))} \\
\texttt{|| ((\colorbox{pink!50}{(\_\_int128)} b == 1) \&\& (\colorbox{pink!50}{(\_\_int128)} i == 1)) || ((\colorbox{pink!50}{(\_\_int128)} b == 1) \&\& (\colorbox{pink!50}{(\_\_int128)} 6 == i)))} \\
\midrule
\textbf{Stage 1: After AST normalization} (99 characters, 66\% reduction): \\
\midrule
\texttt{b == 1 \&\& 5 == i || b == 1 \&\& 2 == i || b == 1 \&\& i == 4 ||} \\
\texttt{b == 1 \&\& 3 == i || b == 1 \&\& i == 1 || b == 1 \&\& 6 == i} \\
\midrule
\textbf{Stage 2: After LLM simplification} (14 characters, \textbf{95\% total reduction}): \\
\midrule
\texttt{1 <= i \&\& i <= 6} \\
\bottomrule
\end{tabular}
\caption{\textbf{Complete pipeline example}. Stage~1 removes \colorbox{pink!50}{12 type casts} and minimizes parentheses (66\% reduction). Stage~2 recognizes that \texttt{b == 1} is always true in context and that the enumerated \texttt{i} values form a contiguous range (additional 86\% reduction). The final invariant is 95\% shorter than the original while remaining sufficient to prove the target property.}
\label{fig:complete-pipeline-example}
\end{figure*}

%% ============================================================================
%% STATISTICS TABLE (generated by scripts/compute_normalization_stats.py)
%% ============================================================================

%% Include the auto-generated statistics table:
%% 
%% Invariant Normalization Statistics
%% Generated by compute_normalization_stats.py

\begin{table}[t]
\centering
\small
\begin{tabular}{@{}lr@{}}
\toprule
\textbf{Transformation} & \textbf{Statistic} \\
\midrule
Total invariants analyzed & 4,000 \\
\midrule
\multicolumn{2}{@{}l}{\textit{Cast Removal}} \\
\quad Invariants affected & 1,792 (44.8\%) \\
\quad Total casts removed & 14,285 \\
\midrule
\multicolumn{2}{@{}l}{\textit{Trivially Decidable Clause Elimination}} \\
\quad Invariants affected & 2,793 (69.8\%) \\
\quad Total clauses eliminated & 13,059 \\
\midrule
\multicolumn{2}{@{}l}{\textit{Parenthesis Minimization}} \\
\quad Total parentheses removed & 138,681 \\
\quad Avg. per invariant & 34.7 \\
\midrule
\textbf{Average length reduction} & \textbf{32.8\%} \\
\bottomrule
\end{tabular}
\caption{Impact of invariant normalization on 4,000 training invariants. Trivially decidable clauses include tautologies ($x \leq x$, $0 < 1$) and contradictions ($1 \leq 0$).}
\label{tab:normalization-stats}
\end{table}


%% Or copy-paste the table directly from normalization_stats_table.tex

%% ============================================================================
%% LLM SIMPLIFICATION STATISTICS TABLE
%% ============================================================================

\begin{table}[t]
\centering
\small
\begin{tabular}{@{}lrr@{}}
\toprule
\textbf{Metric} & \textbf{Stage 1 (AST)} & \textbf{Stage 1+2} \\
\midrule
Invariants with casts & 44.8\% & --- \\
Invariants with tautologies & 69.8\% & --- \\
Verbose after Stage 1 & --- & 72.3\% \\
\midrule
Avg.\ length reduction & 32.8\% & 58.4\% \\
Max length reduction & 78.2\% & 95.1\% \\
\midrule
Quality grade 3 (speedup) & --- & 41.2\% \\
Quality grade 2 (no speedup) & --- & 35.6\% \\
Quality grade 1 (not useful) & --- & 8.4\% \\
Verification failures & 0\% & 14.8\% \\
\bottomrule
\end{tabular}
\caption{Pipeline statistics on 4{,}000 invariants. Stage~1 (AST normalization) achieves 100\% soundness by design. Stage~2 (LLM simplification) produces higher compression but 14.8\% of candidates fail verification; these are filtered out by the verification-in-the-loop.}
\label{tab:pipeline-stats}
\end{table}

%% ============================================================================
%% APPENDIX CONTENT (if needed)
%% ============================================================================

%% \appendix
%% \section{Invariant Normalization Details}
%% \label{app:normalization}
%% 
%% [Include the full definitions, rules, and algorithm here if space-constrained
%% in the main paper]

%% \section{LLM Simplification Details}
%% \label{app:llm-simplification}
%%
%% \subsection{Prompt Template}
%% The system prompt instructs the LLM to:
%% \begin{itemize}
%%     \item Produce logically weaker but still inductive invariants
%%     \item Prefer linear arithmetic (verifiers struggle with non-linear math)
%%     \item Generalize enumerated values to ranges
%%     \item Remove tautological and redundant constraints
%%     \item Use only ASCII characters (no Unicode)
%% \end{itemize}
%%
%% \subsection{Implementation Details}
%% \begin{itemize}
%%     \item Model: Kimi-K2-Thinking (reasoning-capable LLM)
%%     \item Temperature: 1.0 (for diversity in Best-of-N)
%%     \item N: 4 candidates per verbose invariant
%%     \item Verification timeout: 600 seconds per check
%%     \item Parallelism: All N candidates verified concurrently; each candidate's 
%%           correctness and usefulness checks also run in parallel
%% \end{itemize}
